\begin{proposition} \label{proposition:epsilon_delta_definition_of_limit_of_functions_on_extended_metric_spaces_are_equivalent_to_topological_limit_definition}
    Let $T$ be an \CrefAndHyperrefIfExist{definition:ordered_semiring}{ordered semiring}, let $T^\infty$ be the \CrefAndHyperrefIfExist{definition:ordered_semiring_with_adjoined_infinity}{ordered semiring with adjoined infinity} obtained from $T$, and let $(X,d_X)$, $(Y,d_Y)$ be \CrefAndHyperrefIfExist{definition:extended_metric_on_a_set_valued_in_an_ordered_semiring_with_adjoined_infinity}{extended metric spaces valued in $T^\infty$}. Let $f : X \to Y$, let $a$ be an \CrefAndHyperrefIfExist{definition:accumulation_point_in_an_extended_metric_space}{accumulation point} of $X$, and let $L \in Y$. Recall that $X$ and $Y$ have \CrefAndHyperrefIfExist{definition:topology_induced_by_an_extended_metric_and_open_ball_for_an_extended_metric}{induced topological space structures}.

    $f$ has a limit $L$ at $a$ in the sense of \Cref{definition:limit_of_a_function_between_extended_metric_spaces_at_an_accumulation_point} if and only if it does so in the sense of \Cref{definition:limit_of_a_function_between_topological_spaces_at_an_accumulation_point}.
\end{proposition}
\begin{proof}
    Suppose that $\lim_{x \to a} f(x) = L$ in the sense of \Cref{definition:limit_of_a_function_between_extended_metric_spaces_at_an_accumulation_point} holds. Let $V$ be an open neighborhood of $L$ in $Y$. By \Cref{definition:topology_induced_by_an_extended_metric_and_open_ball_for_an_extended_metric}, there exists $\varepsilon \in T_{>0}$ such that the open ball $B(L, \varepsilon) = \{y \in Y : d_Y(y,L) < \varepsilon\}$ satisfies $B(L, \varepsilon) \subseteq V$. By the $\varepsilon$-$\delta$ definition of limit, there exists $\delta \in T_{>0}$ such that for any $x \in X$, the statement $0 < d_X(x,a) < \delta$ implies $d_Y(f(x), L) < \varepsilon$. Equivalently, $x \in B(a, \delta) \setminus \{a\}$ implies $f(x) \in B(L, \varepsilon)$. Thus $f(B(a, \delta) \setminus \{a\}) \subseteq V$, so $\lim_{x \to a} f(x) = L$ in the sense of \Cref{definition:limit_of_a_function_between_topological_spaces_at_an_accumulation_point} holds.

    Conversely, suppose that $\lim_{x \to a} f(x) = L$ in the sense of \Cref{definition:limit_of_a_function_between_topological_spaces_at_an_accumulation_point} holds. Given $\varepsilon \in T_{>0}$, the open ball $B(L, \varepsilon) = \{y \in Y : d_Y(y,L) < \varepsilon\}$ is an open neighborhood of $L$ in $Y$. By the topological definition of limit, there exists an open neighborhood $U$ of $a$ in $X$ such that $f(U \setminus \{a\}) \subseteq B(L, \varepsilon)$. Because $U$ is open and contains $a$, there exists $\delta \in T_{>0}$ such that $B(a, \delta) \subseteq U$. We then have that $x \in B(a, \delta) \setminus \{a\}$ implies $f(x) \in B(L, \varepsilon)$, i.e., $0 < d_X(x,a) < \delta$ implies $d_Y(f(x), L) < \varepsilon$. Therefore, $\lim_{x \to a} f(x) = L$ in the sense of \Cref{definition:limit_of_a_function_between_extended_metric_spaces_at_an_accumulation_point} holds.
\end{proof}